% from ega0-1-fr.tex
\oldpage[0\textsubscript{I}]{41}
\subsection{Image réciproque d'un $\mathcal{B}$-Module}
\label{subsection:0.4.3-fr}

\begin{env}[4.3.1]
\label{0.4.3.1-fr}
Les hypothèses et notations étant celles de (4.2.1), soit $\mathcal{G}$ un
$\mathcal{B}$-Module et $\psi^*(\mathcal{G})$ son image réciproque (3.7.1)
qui est donc un faisceau de groupes abéliens sur $X$. La définition des
sections de $\psi^*(\mathcal{G})$ et de $\psi^*(\mathcal{B})$ (3.7.1)
montre que $\psi^*(\mathcal{G})$ est canoniquement muni d'une structure de
$\psi^*(\mathcal{B})$-Module. D'autre part, l'homomorphisme
$\theta^\sharp:\psi^*(\mathcal{B})\to\mathcal{A}$ munit $\mathcal{A}$ d'une
structure de $\psi^*(\mathcal{B})$-Module, que nous noterons
$\mathcal{A}_{[\theta]}$ lorsqu'il y a lieu d'éviter des confusions~; le
produit tensoriel
$\psi^*(\mathcal{G})\otimes_{\psi^*(\mathcal{B})}\mathcal{A}_{[\theta]}$
est alors muni d'une structure de $\mathcal{A}$-Module. Nous dirons que ce
$\mathcal{A}$-Module est l'image réciproque de $\mathcal{G}$ par le morphisme
$\Psi$ et nous le noterons $\Psi^*(\mathcal{G})$. Si $\mathcal{G}_1$,
$\mathcal{G}_2$ sont deux $\mathcal{B}$-Modules sur $Y$, $v$ un
$\mathcal{B}$-homomorphisme $\mathcal{G}_1\to\mathcal{G}_2$,
$\psi^*(v)$, comme on le vérifie aussitôt, est un
$\psi^*(\mathcal{B})$-homomorphisme de $\psi^*(\mathcal{G}_1)$ dans
$\psi^*(\mathcal{G}_2)$~; par suite $\psi^*(v)\otimes1$ est un
$\mathcal{A}$-homomorphisme $\Psi^*(\mathcal{G}_1)\to\Psi^*(\mathcal{G}_2)$,
que nous noterons $\Psi^*(v)$. On a donc défini $\Psi^*$ comme un
\emph{foncteur covariant} de la catégorie des $\mathcal{B}$-Modules dans celle
des $\mathcal{A}$-Modules. Ici, ce foncteur (contrairement à $\psi^*$) n'est
plus exact en général, mais seulement \emph{exact à droite}, la tensorisation
par $\mathcal{A}$ étant un foncteur exact à droite dans la catégorie des
$\psi^*(\mathcal{B})$-Modules.

Pour tout $x\in X$, on a
\[
  (\Psi^*(\mathcal{G}))_x
  =\mathcal{G}_{\psi(x)}\otimes_{\mathcal{B}_{\psi(x)}}\mathcal{A}_x,
\]
en vertu de (3.7.2). Le support de $\Psi^*(\mathcal{G})$ est donc contenu dans
$\psi^{-1}(\operatorname{Supp}\mathcal{G})$.
\end{env}

\begin{env}[4.3.2]
\label{0.4.3.2-fr}
Soit $(\mathcal{G}_\lambda)$ un système inductif de $\mathcal{B}$-Modules, et
soit $\mathcal{G}=\varinjlim\mathcal{G}_\lambda$ sa limite inductive. Les
homomorphismes canoniques $\mathcal{G}_\lambda\to\mathcal{G}$ définissent des
$\psi^*(\mathcal{B})$-homomorphismes
$\psi^*(\mathcal{G}_\lambda)\to\psi^*(\mathcal{G})$, qui donnent un
homomorphisme canonique
\[
  \varinjlim\psi^*(\mathcal{G}_\lambda)
  \longrightarrow
  \psi^*(\mathcal{G}).
\]
Comme la fibre en un point d'une limite inductive de faisceaux est la limite
inductive des fibres au même point (G, II, 1.11), l'homomorphisme canonique
précédent est \emph{bijectif} (3.7.2). En outre, le produit tensoriel commute
aux limites inductives de faisceaux, et on a donc un \emph{isomorphisme
canonique fonctoriel}
\[
  \varinjlim\Psi^*(\mathcal{G}_\lambda)
  \xrightarrow{\sim}
  \Psi^*\bigl(\varinjlim\mathcal{G}_\lambda\bigr)
\]
de $\mathcal{A}$-Modules.

D'autre part, pour une somme directe finie $\bigoplus_i\mathcal{G}_i$ de
$\mathcal{B}$-Modules, il est clair que
\[
  \psi^*\bigl(\bigoplus_i\mathcal{G}_i\bigr)
  =\bigoplus_i\psi^*(\mathcal{G}_i)~;
\]
donc, par produit tensoriel avec $\mathcal{A}_{[\theta]}$,
\[
  \Psi^*\bigl(\bigoplus_i\mathcal{G}_i\bigr)
  =\bigoplus_i\Psi^*(\mathcal{G}_i).
  \tag{4.3.2.1}\label{0.4.3.2.1-fr}
\]
Par passage à la limite inductive, on en déduit, en vertu de ce qui précède,
que l'égalité précédente est encore vraie pour une somme directe
\emph{quelconque}.
\end{env}

\begin{env}[4.3.3]
\EGAReaderTarget{canon-ref.0.4.3.3-fr}{0.4.3.3-fr}
Soient $\mathcal{G}_1$, $\mathcal{G}_2$ deux $\mathcal{B}$-Modules~; de la
définition des images réciproques de faisceaux de groupes abéliens (3.7.1),
on déduit aussitôt un homomorphisme canonique
\[
  \psi^*(\mathcal{G}_1)\otimes_{\psi^*(\mathcal{B})}\psi^*(\mathcal{G}_2)
  \longrightarrow
  \psi^*(\mathcal{G}_1\otimes_{\mathcal{B}}\mathcal{G}_2)
\]
de $\psi^*(\mathcal{B})$-Modules, et la fibre en un point d'un produit
tensoriel de faisceaux étant le produit tensoriel des fibres en ce point
(G, II, 2.8), on déduit de (3.7.2) que l'homomorphisme précédent est en fait
un \emph{isomorphisme}. Par produit tensoriel avec $\mathcal{A}$, on en déduit
un \emph{isomorphisme canonique fonctoriel}
\[
  \Psi^*(\mathcal{G}_1)\otimes_{\mathcal{A}}\Psi^*(\mathcal{G}_2)
  \xrightarrow{\sim}
  \Psi^*(\mathcal{G}_1\otimes_{\mathcal{B}}\mathcal{G}_2).
  \tag{4.3.3.1}\label{0.4.3.3.1-fr}
\]
\end{env}

\begin{env}[4.3.4]
\label{0.4.3.4-fr}
Soit $\mathcal{C}$ une $\mathcal{B}$-Algèbre~; la donnée de la structure
d'Algèbre sur $\mathcal{C}$ revient à la donnée d'un $\mathcal{B}$-homomorphisme
\[
  \mathcal{C}\otimes_{\mathcal{B}}\mathcal{C}\longrightarrow\mathcal{C}
\]
satisfaisant aux conditions d'associativité et de commutativité (conditions
qui se vérifient sur chaque fibre)~; l'isomorphisme précédent permet de
transporter cet homomorphisme en un homomorphisme de $\mathcal{A}$-Modules
\[
  \Psi^*(\mathcal{C})\otimes_{\mathcal{A}}\Psi^*(\mathcal{C})
  \longrightarrow\Psi^*(\mathcal{C})
\]
satisfaisant aux mêmes conditions, donc $\Psi^*(\mathcal{C})$ est ainsi muni
d'une structure de $\mathcal{A}$-Algèbre. En particulier, il résulte aussitôt
des définitions que la $\mathcal{A}$-Algèbre $\Psi^*(\mathcal{B})$ est égale à
$\mathcal{A}$ (à un isomorphisme canonique près).

De même, si $\mathcal{M}$ est un $\mathcal{C}$-Module, la donnée de cette
structure de Module revient
